\documentclass[11pt,a4paper]{article}

\usepackage[margin=2.3cm]{geometry}
\usepackage{amsmath,amssymb}
\usepackage{booktabs}
\usepackage{longtable}
\usepackage{array}
\usepackage{float}
\usepackage{graphicx}
\usepackage{hyperref}
\usepackage{listings}
\usepackage{microtype}

\hypersetup{
    colorlinks=true,
    linkcolor=blue,
    urlcolor=blue,
    citecolor=blue
}

\lstset{
    basicstyle=\ttfamily\small,
    frame=single,
    breaklines=true,
    columns=fullflexible
}

\title{\textbf{Report on Text Similarity and Tree Estimation Exercise}\\

\large DA7067 Computational Mathematics}

\author{ Onyinyechi Sonia Bäckström, Yuqi Yang, Lovisa Irberger\\
Stockholm University}

\date{September 2026}

\begin{document}

\maketitle
\newpage
\tableofcontents
\newpage
\begin{abstract}

This report describes a Python implementation for estimating genomic similarity
using k-mers, density-based MinHash sketching, and the Jaccard Index. The
estimated Jaccard values are converted to genomic distances using
$d=1-J$. These distances are then used to construct binary evolutionary trees
with UPGMA.

The program accepts a FASTA file, a k-mer size, and a sketching density as
command-line parameters. Experiments were performed on three simulated
datasets containing 3, 5, and 10 genomes, and on a real dataset containing
20 \textit{E. coli} genomes.

The report investigates how sketching density and k-mer size affect the
estimated distance, evaluates the resulting evolutionary trees against the
known lineage information in the simulated datasets, and reports the runtime
of the implementation. Exact Jaccard values and distances are also calculated
for the small simulated datasets as an additional validation of the sketch
estimates.

\end{abstract}
\newpage


% ============================================================
\section{Introduction}
% ============================================================

The purpose of this assignment is to estimate pairwise genomic distances using
k-mer sketching and the Jaccard Index, and then use the resulting distances to
construct a binary evolutionary tree.

For two genomes $A$ and $B$, the genomic distance is defined as

\begin{equation}
d(A,B)=1-J(A,B),
\end{equation}

where $J(A,B)$ is the Jaccard Index.

Task 1 investigates two questions:

\begin{enumerate}
    \item Does the estimated genomic distance change with sketching density?
    \item Does the estimated genomic distance change with k-mer size?
\end{enumerate}

Task 2 uses the estimated distances from Task 1 to construct binary
evolutionary trees. The assignment requires the tree structure, rather than
branch lengths. The three simulated datasets contain 3, 5, and 10 genomes,
while the real dataset contains 20 \textit{E. coli} genomes.

For the simulated datasets, the accession numbers also contain lineage
information. This makes it possible to compare the estimated trees with the
known simulated lineage.


% ============================================================
\section{Program and How to Run It}
% ============================================================

\subsection{Programming Language}

The program is written in Python 3.

The implementation uses the Python standard-library modules
\texttt{argparse}, \texttt{hashlib}, \texttt{time}, and \texttt{csv}.

\subsection{Command-Line Usage}

The program is run using

\begin{lstlisting}
python Problem3.py <FASTA file> <k-mer size> <density>
\end{lstlisting}

For example:

\begin{lstlisting}
python Problem3.py test1.fa 30 0.5
python Problem3.py test2.fa 30 0.5
python Problem3.py test3.fa 30 0.5
python Problem3.py ecoli_20_genomes_short_acc.fa 30 0.5
\end{lstlisting}

The program reads the genomes from the FASTA file and reports the number of
k-mers, sketch size, pairwise estimated Jaccard values, estimated distances,
the distance matrix, the UPGMA tree, and runtime information.

The estimated distance matrix is also saved to

\begin{lstlisting}
distances.csv
\end{lstlisting}


% ============================================================
\section{Methods}
% ============================================================

\subsection{Reading FASTA Files}

Each genome is read from the FASTA file. A line beginning with \texttt{>}
indicates a new genome, and the following sequence lines are joined together.

For a sequence of length $L$ and k-mer size $k$, the program considers

\begin{equation}
L-k+1
\end{equation}

possible k-mers.

Only k-mers containing the characters A, C, G, or T are retained. The k-mers
are stored in a Python set, meaning that repeated copies of the same k-mer
are represented only once.

For the density experiment, $k=30$ was used with densities
$0.1, 0.2, 0.3, 0.5, 0.7,$ and $0.9$. For the k-mer-size experiment,
density $=0.5$ was used with $k=10, 30, 50, 70,$ and $90$. The main
Task 2 experiments and the real \textit{E. coli} experiment used
$k=30$ and density $=0.5$.
\subsection{Hashing and Sketching}

Each k-mer is converted into a deterministic 64-bit hash using BLAKE2b.
The hash value is normalized by

\begin{equation}
h_{\mathrm{norm}}=\frac{h}{2^{64}}.
\end{equation}

A hash is included in the sketch if

\begin{equation}
h_{\mathrm{norm}} < p,
\end{equation}

where $p$ is the user-specified sketching density.

For example, a density of $0.5$ keeps approximately half of the hash values,
while a density of $0.1$ produces a much smaller sketch.

This is a density-based hash sketching approach from the MinHash methods
discussed in the lectures.

A deterministic BLAKE2b hash was used instead of Python's built-in
\texttt{hash()} because the built-in hash can use a different randomized seed
between program runs.

\subsection{Jaccard Similarity and Genomic Distance}

For two sketches $S_A$ and $S_B$, the estimated Jaccard Index is

\begin{equation}
\widehat{J}(A,B)
=
\frac{|S_A\cap S_B|}
{|S_A\cup S_B|}.
\end{equation}

The estimated genomic distance is then

\begin{equation}
\widehat{d}(A,B)
=
1-\widehat{J}(A,B).
\end{equation}

The program calculates this for every pair of genomes and stores the results
in a symmetric distance matrix.

\subsection{UPGMA Tree Construction}

The distance matrix is used to construct a binary tree with UPGMA
(Unweighted Pair Group Method with Arithmetic Mean).

Initially, every genome is a separate cluster. The two closest clusters are
merged repeatedly until only one cluster remains.

When clusters $A$ and $B$ are merged, their distance to another cluster $C$
is calculated using

\begin{equation}
d(A\cup B,C)
=
\frac{|A|d(A,C)+|B|d(B,C)}
{|A|+|B|}.
\end{equation}

Here, $|A|$ and $|B|$ are the numbers of original genomes in the two
clusters.

The final result is a binary tree containing the original genomes as leaves.
Only the structure of the tree is considered in this assignment.

\subsection{Accession Numbers and Genome Labels}

An accession is the genome identifier provided in the FASTA file.

For the simulated datasets, the accession number also contains information
about the intended lineage. The program renames the genomes as G1, G2, G3,
etc. to make the output easier to read.

For example, in Test 1:

\begin{itemize}
    \item G1 = accession 1, lineage [0]
    \item G2 = accession 5, lineage [2,0]
    \item G3 = accession 6, lineage [2,0]
\end{itemize}

Thus, G2 and G3 have the same simulated lineage.

\subsection{Runtime Measurement}

The program measures four parts of the computation:

\begin{enumerate}
    \item sketch generation,
    \item distance calculation,
    \item UPGMA construction,
    \item total computation time.
\end{enumerate}

The measurements are obtained using Python's
\texttt{time.perf\_counter()}.


% ============================================================
\section{Task 1 Results}
% ============================================================

\subsection{Effect of Sketching Density}

To investigate the first guiding question, Test 1 was run with a fixed
k-mer size of $k=30$. The densities tested were
\[
0.1,\ 0.2,\ 0.3,\ 0.5,\ 0.7,\ 0.9.
\]

The main Task 1 result is the estimated Jaccard Index and the corresponding
estimated genomic distance.
The number of distinct k-mers was 171 for each of the three genomes at
every density, since the k-mer size was kept fixed. The sketch sizes
increased as the density increased. At density 0.1, the sketch sizes of
G1, G2, and G3 were 25, 16, and 19, respectively. At density 0.2 they
were 40, 35, and 32; at density 0.3 they were 57, 59, and 54; at density
0.5 they were 88, 92, and 92; at density 0.7 they were 120, 123, and
123; and at density 0.9 they were 149, 161, and 155.

The estimated distance changes with sketching density. For example, for
G2--G3 the estimated distance changes from 0.6000 at density 0.1 to
0.6610 at density 0.9.

The change is not perfectly monotonic for every genome pair. However,
increasing the density produces a larger sketch and retains more
information from the original k-mer sets.

Therefore, the answer to the first guiding question is:
\begin{quote}
\textbf{The estimated genomic distance changes with sketching density.}
\end{quote}
\begin{table}[H]
\centering
\small
\caption{Estimated Jaccard values and genomic distances for different
sketching densities on Test 1 with $k=30$.}
\begin{tabular}{c|c|c|c}
\toprule
Density & Genome pair & Estimated $J$ & Estimated $d$\\
\midrule
0.1 & G1--G2 & 0.0789 & 0.9211\\
    & G1--G3 & 0.0732 & 0.9268\\
    & G2--G3 & 0.4000 & 0.6000\\
\midrule
0.2 & G1--G2 & 0.0714 & 0.9286\\
    & G1--G3 & 0.0746 & 0.9254\\
    & G2--G3 & 0.4255 & 0.5745\\
\midrule
0.3 & G1--G2 & 0.0841 & 0.9159\\
    & G1--G3 & 0.0882 & 0.9118\\
    & G2--G3 & 0.4125 & 0.5875\\
\midrule
0.5 & G1--G2 & 0.0843 & 0.9157\\
    & G1--G3 & 0.0843 & 0.9157\\
    & G2--G3 & 0.3731 & 0.6269\\
\midrule
0.7 & G1--G2 & 0.0897 & 0.9103\\
    & G1--G3 & 0.0897 & 0.9103\\
    & G2--G3 & 0.3443 & 0.6557\\
\midrule
0.9 & G1--G2 & 0.0839 & 0.9161\\
    & G1--G3 & 0.0857 & 0.9143\\
    & G2--G3 & 0.3390 & 0.6610\\
\bottomrule
\end{tabular}
\end{table}


% ============================================================
\subsection{Effect of K-mer Size}
% ============================================================
To investigate the second guiding question, Test 1 was run with a fixed sketching density of $0.5$. The tested k-mer sizes were
$10$, $30$, $50$, $70$, and $90$.

The sketch sizes for G1, G2, and G3 were 104, 105, and 102 at $k=10$.
At $k=30$, the sketch sizes were 88, 92, and 92. At $k=50$, they were
69, 77, and 76. At $k=70$, they were 66, 79, and 69, and at $k=90$ they
were 61, 56, and 58.

For $k=30$, each genome contained 171 distinct k-mers. 

The estimated Jaccard values and genomic distances were:

\begin{table}[H]
\centering
\small
\caption{Estimated Jaccard values and genomic distances for different
k-mer sizes on Test 1 with density $=0.5$.}
\begin{tabular}{c|c|c|c}
\toprule
$k$ & Genome pair & Estimated $J$ & Estimated $d$\\
\midrule
10 & G1--G2 & 0.4027 & 0.5973\\
   & G1--G3 & 0.3826 & 0.6174\\
   & G2--G3 & 0.5682 & 0.4318\\
\midrule
30 & G1--G2 & 0.0843 & 0.9157\\
   & G1--G3 & 0.0843 & 0.9157\\
   & G2--G3 & 0.3731 & 0.6269\\
\midrule
50 & G1--G2 & 0.0000 & 1.0000\\
   & G1--G3 & 0.0000 & 1.0000\\
   & G2--G3 & 0.2750 & 0.7250\\
\midrule
70 & G1--G2 & 0.0000 & 1.0000\\
   & G1--G3 & 0.0000 & 1.0000\\
   & G2--G3 & 0.2333 & 0.7667\\
\midrule
90 & G1--G2 & 0.0000 & 1.0000\\
   & G1--G3 & 0.0000 & 1.0000\\
   & G2--G3 & 0.1400 & 0.8600\\
\bottomrule
\end{tabular}
\end{table}

The k-mer size has a strong effect on the estimated genomic distance.
For G2--G3, the estimated distance increases from 0.4318 at $k=10$ to
0.8600 at $k=90$.

For G1--G2 and G1--G3, the estimated Jaccard Index is zero for
$k=50$, $70$, and $90$, resulting in an estimated distance of 1.0000.

Therefore, the answer to the second guiding question is:

\begin{quote}
\textbf{The genomic distance changes substantially with k-mer size.}
\end{quote}

Larger k-mers are more specific, which results in fewer shared k-mers
between the genomes in these experiments.

% ============================================================
\section{Additional Validation}
% ============================================================

The exact Jaccard Index and exact genomic distance are not required by the
assignment. They are included here as a separate validation of the sketch
estimates.

For the small simulated datasets, the complete k-mer sets can be used to
calculate

\begin{equation}
J_{\mathrm{exact}}(A,B)
=
\frac{|K_A\cap K_B|}
{|K_A\cup K_B|}.
\end{equation}

The exact distance is

\begin{equation}
d_{\mathrm{exact}}
=
1-J_{\mathrm{exact}}.
\end{equation}

The error is

\begin{equation}
\mathrm{Error}
=
d_{\mathrm{estimated}}-d_{\mathrm{exact}},
\end{equation}

and the absolute error is

\begin{equation}
\mathrm{Abs.Error}
=
|d_{\mathrm{estimated}}-d_{\mathrm{exact}}|.
\end{equation}

These values are only used to check how close the sketch estimates are to the
values obtained from the complete k-mer sets. They are not used instead of the
estimated distances in the main analysis or in the UPGMA tree construction.

\subsection{Validation: Sketching Density}

\begin{longtable}{llrrrrrr}
\caption{Exact and estimated values for different sketching densities on
Test 1 with $k=30$.}
\\
\toprule
Density & Pair & Exact $J$ & Est. $J$ &
Exact $d$ & Est. $d$ & Error & Abs. Error\\
\midrule
\endfirsthead

\toprule
Density & Pair & Exact $J$ & Est. $J$ &
Exact $d$ & Est. $d$ & Error & Abs. Error\\
\midrule
\endhead

0.1&G1--G2&0.0755&0.0789&0.9245&0.9211&-0.0035&0.0035\\
0.1&G1--G3&0.0755&0.0732&0.9245&0.9268&0.0023&0.0023\\
0.1&G2--G3&0.3307&0.4000&0.6693&0.6000&-0.0693&0.0693\\

0.2&G1--G2&0.0755&0.0714&0.9245&0.9286&0.0040&0.0040\\
0.2&G1--G3&0.0755&0.0746&0.9245&0.9254&0.0008&0.0008\\
0.2&G2--G3&0.3307&0.4255&0.6693&0.5745&-0.0948&0.0948\\

0.3&G1--G2&0.0755&0.0841&0.9245&0.9159&-0.0086&0.0086\\
0.3&G1--G3&0.0755&0.0882&0.9245&0.9118&-0.0128&0.0128\\
0.3&G2--G3&0.3307&0.4125&0.6693&0.5875&-0.0818&0.0818\\

0.5&G1--G2&0.0755&0.0843&0.9245&0.9157&-0.0089&0.0089\\
0.5&G1--G3&0.0755&0.0843&0.9245&0.9157&-0.0089&0.0089\\
0.5&G2--G3&0.3307&0.3731&0.6693&0.6269&-0.0424&0.0424\\

0.7&G1--G2&0.0755&0.0897&0.9245&0.9103&-0.0142&0.0142\\
0.7&G1--G3&0.0755&0.0897&0.9245&0.9103&-0.0142&0.0142\\
0.7&G2--G3&0.3307&0.3443&0.6693&0.6557&-0.0135&0.0135\\

0.9&G1--G2&0.0755&0.0839&0.9245&0.9161&-0.0084&0.0084\\
0.9&G1--G3&0.0755&0.0857&0.9245&0.9143&-0.0102&0.0102\\
0.9&G2--G3&0.3307&0.3390&0.6693&0.6610&-0.0082&0.0082\\

\bottomrule
\end{longtable}

The exact values remain constant because changing the density does not change
the original k-mer sets. Only the sketch changes.

For G2--G3, the absolute error was 0.0693 at density 0.1 and 0.0082 at
density 0.9. The error was not perfectly monotonic, however. For example, it
increased from 0.0693 to 0.0948 when the density increased from 0.1 to 0.2.

Thus, higher density generally provided a more accurate estimate in this
experiment, but it did not guarantee a smaller error at every density.


% ============================================================
\subsection{Validation: K-mer Size}
% ============================================================

\begin{longtable}{llrrrrrr}
\caption{Exact and estimated values for different k-mer sizes on Test 1
with density 0.5.}
\\
\toprule
$k$ & Pair & Exact $J$ & Est. $J$ &
Exact $d$ & Est. $d$ & Error & Abs. Error\\
\midrule
\endfirsthead

\toprule
$k$ & Pair & Exact $J$ & Est. $J$ &
Exact $d$ & Est. $d$ & Error & Abs. Error\\
\midrule
\endhead

10&G1--G2&0.3791&0.4027&0.6209&0.5973&-0.0236&0.0236\\
10&G1--G3&0.3357&0.3826&0.6643&0.6174&-0.0469&0.0469\\
10&G2--G3&0.5159&0.5682&0.4841&0.4318&-0.0523&0.0523\\

30&G1--G2&0.0755&0.0843&0.9245&0.9157&-0.0089&0.0089\\
30&G1--G3&0.0755&0.0843&0.9245&0.9157&-0.0089&0.0089\\
30&G2--G3&0.3307&0.3731&0.6693&0.6269&-0.0424&0.0424\\

50&G1--G2&0.0000&0.0000&1.0000&1.0000&0.0000&0.0000\\
50&G1--G3&0.0000&0.0000&1.0000&1.0000&0.0000&0.0000\\
50&G2--G3&0.2743&0.2750&0.7257&0.7250&-0.0007&0.0007\\

70&G1--G2&0.0000&0.0000&1.0000&1.0000&0.0000&0.0000\\
70&G1--G3&0.0000&0.0000&1.0000&1.0000&0.0000&0.0000\\
70&G2--G3&0.2074&0.2333&0.7926&0.7667&-0.0260&0.0260\\

90&G1--G2&0.0000&0.0000&1.0000&1.0000&0.0000&0.0000\\
90&G1--G3&0.0000&0.0000&1.0000&1.0000&0.0000&0.0000\\
90&G2--G3&0.1269&0.1400&0.8731&0.8600&-0.0131&0.0131\\

\bottomrule
\end{longtable}

The validation shows the same general pattern as the main Task 1 results.

For G2--G3, the exact Jaccard Index decreased from 0.5159 at $k=10$ to
0.1269 at $k=90$. Consequently, the exact distance increased from 0.4841
to 0.8731.

The sketch estimates were close to the exact values in several cases. For
example, at $k=50$, the absolute error for G2--G3 was only 0.0007.

The validation therefore supports the conclusion that the choice of k-mer size
has a substantial effect on the measured similarity.


% ============================================================
\section{Task 2: Evolutionary Tree Estimation}
% ============================================================

All three simulated datasets were run using

\[
k=30,\qquad p=0.5.
\]

The trees were constructed using the estimated distances from the sketches.

\subsection{Test 1 Dataset: Three Genomes}

The lineage information was:

\begin{itemize}
    \item G1 = accession 1, lineage [0]
    \item G2 = accession 5, lineage [2,0]
    \item G3 = accession 6, lineage [2,0]
\end{itemize}

The estimated pairwise distances were:

\begin{itemize}
    \item G1--G2: 0.9157
    \item G1--G3: 0.9157
    \item G2--G3: 0.6269
\end{itemize}

The resulting tree was:

\begin{lstlisting}
Node2
  G1
  Node1
    G2
    G3
\end{lstlisting}

G2 and G3 have the smallest estimated distance and are grouped together.
They also have the same lineage [2,0].

Therefore, the tree agrees with the simulated lineage for Test 1.


% ------------------------------------------------------------
\subsection{Test 2 Dataset: Five Genomes}

The lineage information was:

\begin{itemize}
    \item G1 = accession 2, lineage [0]
    \item G2 = accession 7, lineage [3,1,0]
    \item G3 = accession 8, lineage [3,1,0]
    \item G4 = accession 9, lineage [4,1,0]
    \item G5 = accession 10, lineage [4,1,0]
\end{itemize}

The estimated distances were mostly $1.0000$, indicating very low estimated Jaccard similarity between most genome pairs. The smallest distance was between G4 and G5 ($0.9739$), followed by G2 and G4 ($0.9876$).

The estimated and exact pairwise distances  were:
\begin{table}[h]
\centering
\begin{tabular}{lcc}
\hline
\textbf{Pair} & \textbf{Exact distance} & \textbf{Estimated distance} \\
\hline
G1--G2 & 1.0000 & 1.0000 \\
G1--G3 & 1.0000 & 1.0000 \\
G1--G4 & 1.0000 & 1.0000 \\
G1--G5 & 1.0000 & 1.0000 \\
G2--G3 & 0.9971 & 1.0000 \\
G2--G4 & 0.9882 & 0.9876 \\
G2--G5 & 1.0000 & 1.0000 \\
G3--G4 & 1.0000 & 1.0000 \\
G3--G5 & 1.0000 & 1.0000 \\
G4--G5 & 0.9760 & 0.9739 \\
\hline
\end{tabular}
\caption{Exact and estimated pairwise distances for Test 2 using
$k=30$ and density $=0.5$.}
\label{tab:test2_distances}
\end{table}

The resulting tree was:

\begin{lstlisting}
Node4
  Node2
    G2
    Node1
      G4
      G5
  Node3
    G1
    G3
\end{lstlisting}

G4 and G5 have the same lineage, have the smallest estimated distance and are grouped directly together in the estimated tree. 
G2 and G3 also have the same lineage, but they are not grouped directly together. 
Instead, the estimated tree groups G1 and G3 together. 
The exact distance between G2 and G3 is $0.9971$, while their estimated distance is $1.0000$. 
Therefore, the estimated tree does not fully agree with the simulated lineage structure in Test 2 Dataset, although the G4--G5 pair is correctly grouped.



% ------------------------------------------------------------
\subsection{Test 3 Dataset: Ten Genomes}
For Test 3 Dataset, we used $k=30$ and a density of $0.5$. Each genome contained
approximately 1,971 distinct k-mers. The resulting sketch sizes were
975, 1004, 965, 984, 976, 951, 1007, 1030, 991, and 988 for G1--G10,
respectively.


The lineage information was:

\begin{itemize}
    \item G1 = accession 4, lineage [1,0]
    \item G2 = accession 5, lineage [2,0]
    \item G3 = accession 7, lineage [3,1,0]
    \item G4 = accession 17, lineage [8,3,1,0]
    \item G5 = accession 18, lineage [8,3,1,0]
    \item G6 = accession 29, lineage [14,6,2,0]
    \item G7 = accession 27, lineage [13,6,2,0]
    \item G8 = accession 28, lineage [13,6,2,0]
    \item G9 = accession 61, lineage [30,14,6,2,0]
    \item G10 = accession 62, lineage [30,14,6,2,0]
\end{itemize}

The estimated distances for same-lineage pairs are:

\[
\begin{array}{c|c|c}
\text{Genome pair} & \text{Exact distance} & \text{Estimated distance} \\
\hline
G4-G5 & 0.7818 & 0.7841 \\
G7-G8 & 0.8306 & 0.8246 \\
G9-G10 & 0.8134 & 0.8071
\end{array}
\]


The resulting tree was:

\begin{lstlisting}
Node9
  Node6
    G1
    Node5
      G3
      Node1
        G4
        G5
  Node8
    Node4
      Node2
        G9
        G10
      Node7
        G2
        Node3
          G7
          G8
\end{lstlisting}

The estimated tree showed good agreement with the simulated lineage
structure. The genomes G4 and G5, G7 and G8, and G9 and G10 have the same
lineage and were grouped directly together in the estimated tree. Their
estimated distances were $0.7841$, $0.8246$, and $0.8071$, respectively.
The corresponding exact distances were $0.7818$, $0.8306$, and $0.8134$.

The tree also reflected some of the deeper relationships between the
genomes. For example, G3 was grouped with the G4--G5 pair, while G6 was
relatively close to G7--G8. Overall, Test 3 Dataset showed stronger agreement
between the estimated tree and the simulated lineage structure than
Test 2 Dataset.


% ------------------------------------------------------------
\subsection{Task 2 Summary}

\begin{table}[H]
\centering
\small
\caption{Comparison of the estimated trees with the simulated lineage
information.}
\begin{tabular}{p{2.2cm}p{4.5cm}p{7cm}}
\toprule
Dataset & Same-lineage pairs & Observation\\
\midrule
Test 1
&
G2--G3
&
G2 and G3 are grouped directly together.
\\
\midrule
Test 2
&
G2--G3; G4--G5
&
G4 and G5 are grouped directly together, but G2 and G3 are not.
\\
\midrule
Test 3
&
G4--G5; G7--G8; G9--G10
&
All three same-lineage pairs are grouped directly together.
\\
\bottomrule
\end{tabular}
\end{table}

These results show that genomes with smaller estimated distances tend to be
placed closer together by UPGMA. However, the tree can be affected by errors
in the estimated distances. Test 2 is an example where two genomes with the
same simulated lineage were not directly grouped.


% ============================================================
\subsection{The Real \textit{E. coli} Dataset}
% ============================================================

The real dataset contains 20 \textit{E. coli} genomes.

The experiment used

\[
k=30,\qquad p=0.5.
\]

Exact Jaccard validation was not performed for this dataset because each
genome contains millions of k-mers.

\begin{longtable}{llrr}
\caption{Genome and sketch sizes for the 20 \textit{E. coli} genomes.}
\\
\toprule
Genome & Accession & Number of k-mers & Sketch size\\
\midrule
\endfirsthead

\toprule
Genome & Accession & Number of k-mers & Sketch size\\
\midrule
\endhead

G1&NZ\_CP116919.1&5,192,305&2,595,531\\
G2&NZ\_CP117068.1&5,136,027&2,549,138\\
G3&NZ\_CP117013.1&5,184,086&2,592,634\\
G4&NZ\_CP117116.1&5,168,088&2,584,664\\
G5&NZ\_CP132204.1&5,122,880&2,512,046\\
G6&NZ\_AP027159.1&5,437,795&2,720,323\\
G7&NZ\_AP027162.1&5,436,737&2,719,778\\
G8&NZ\_AP027165.1&5,322,450&2,665,594\\
G9&NZ\_AP027170.1&5,373,778&2,687,963\\
G10&NZ\_AP027176.1&5,406,046&2,724,746\\
G11&NZ\_AP027181.1&5,477,895&2,740,794\\
G12&NZ\_AP027183.1&5,516,698&2,746,865\\
G13&NZ\_AP027188.1&5,486,795&2,744,471\\
G14&NZ\_AP027191.1&5,517,653&2,768,773\\
G15&NZ\_CP117067.1&5,422,584&2,713,152\\
G16&NZ\_CP116807.1&5,212,812&2,687,365\\
G17&NZ\_CP116451.1&5,162,977&2,511,494\\
G18&NZ\_CP116835.1&4,943,781&2,404,410\\
G19&NZ\_CP123548.1&4,997,126&2,498,412\\
G20&NZ\_CP123549.1&5,016,355&2,508,301\\

\bottomrule
\end{longtable}

The genomes contain approximately 4.9--5.5 million k-mers each, while the
sketches contain approximately 2.4--2.8 million hashes each.

Some of the smallest estimated distances were:

\begin{table}[H]
\centering
\caption{Examples of small estimated distances in the \textit{E. coli}
dataset.}
\begin{tabular}{lrr}
\toprule
Pair & Estimated $J$ & Estimated $d$\\
\midrule
G8--G9   & 0.9561 & 0.0439\\
G8--G10  & 0.9546 & 0.0454\\
G9--G10  & 0.9474 & 0.0526\\
G14--G15 & 0.9463 & 0.0537\\
G11--G15 & 0.9437 & 0.0563\\
G11--G14 & 0.9371 & 0.0629\\
G13--G15 & 0.9369 & 0.0631\\
\bottomrule
\end{tabular}
\end{table}

These results show that some genome pairs are much more similar according to
the k-mer sketches than others.

The estimated distances were used as input to UPGMA. The resulting
evolutionary tree was:

\newpage
\begin{lstlisting}
Node19
  Node17
    Node12
      G17
      G19
    Node14
      G3
      Node10
        G2
        Node9
          G1
          G5
  Node18
    Node11
      G6
      Node8
        Node2
          G10
          Node1
            G8
            G9
        Node7
          G12
          Node6
            G7
            Node5
              G13
              Node4
                G11
                Node3
                  G14
                  G15
    Node16
      G18
      Node15
        G16
        Node13
          G4
          G20
\end{lstlisting}

Several pairs with small estimated distances were grouped directly in
the tree. For example, G8 and G9 form a pair, as do G14 and G15. These
pairs also had some of the smallest estimated distances in the distance matrix. Other directly grouped pairs include G17--G19, G1--G5, G4--G20, and G16 with the G4--G20 group.

Unlike the simulated datasets, the real \textit{E. coli} dataset does
not provide the intended lineage information. 

The total computation time was 449.322 seconds, approximately
7 minutes and 29 seconds. Sketch generation took 328.382 seconds,
distance calculation took 121.011 seconds, and UPGMA took only
0.009 seconds.


% ============================================================
\section{Runtime and Computational Complexity}
% ============================================================

The runtime measurements for the main Task 2 experiments were:

\begin{table}[H]
\centering
\small
\caption{Runtime measurements for the main Task 2 runs.}
\begin{tabular}{lrrrr}
\toprule
Dataset & Sketching & Distance & UPGMA & Total\\
\midrule
Test 1 (3 genomes)
& 0.002 s & 0.000 s & 0.002 s & 0.005 s\\

Test 2 (5 genomes)
& 0.002 s & 0.001 s & 0.002 s & 0.005 s\\

Test 3 (10 genomes)
& 0.054 s & 0.011 s & 0.003 s & 0.068 s\\

\textit{E. coli} (20 genomes)
& 328.382 s & 121.011 s & 0.009 s & 449.322 s\\

\bottomrule
\end{tabular}
\end{table}

The three simulated datasets all took less than one second.

The real \textit{E. coli} dataset took 449.322 seconds, which is approximately
7 minutes and 29 seconds.

The largest part of the runtime was sketch generation, followed by distance
calculation. UPGMA was very fast in comparison.

For $n$ genomes, the number of unordered genome pairs is

\begin{equation}
\frac{n(n-1)}{2}=O(n^2).
\end{equation}

Therefore, the number of pairwise comparisons increases quadratically with
the number of genomes.

For each genome, k-mer generation and hashing depend approximately linearly
on the sequence length. The simple UPGMA implementation repeatedly searches
the current clusters for the closest pair, giving it approximately cubic
behaviour with respect to the number of genomes.

In the experiments, however, UPGMA was not the computational bottleneck.
The main cost came from processing the large genome sequences and their
millions of k-mers.


% ============================================================
\section{Issues, Limitations, and Possible Improvements}
% ============================================================

\subsection{Hash Stability}

Python's built-in \texttt{hash()} can produce different results between
program runs because of hash randomization.

This was avoided by using deterministic BLAKE2b hashing.

\subsection{Sketching Density}

A low density produces a small sketch and therefore requires less data to
compare. However, the density experiment showed that the estimated distance
can change when the sketch is very small.

A possible improvement would be to select the density based on a target sketch
size or memory limit.

\subsection{K-mer Size}

The k-mer size has a strong effect on the measured similarity.

In our experiment, G1--G2 and G1--G3 had zero exact Jaccard similarity for
$k\geq50$.

Therefore, choosing a very large k can make genomes appear completely
different even when they share some shorter sequence patterns.

\subsection{Exact Validation}

Exact Jaccard calculation is practical for the small simulated datasets, but
it becomes much more expensive for the real genomes because they contain
millions of k-mers.

For this reason, exact validation was only performed on the small datasets.

\subsection{Effect of Sketching Error on the Tree}

UPGMA constructs the tree using the estimated distance matrix. Therefore,
differences between the estimated and exact pairwise distances can affect
the resulting tree structure. Test 2 illustrates this effect: G2 and G3
have the same simulated lineage, but their exact distance was $0.9971$
while their estimated distance was $1.0000$. As a result, they were not
directly grouped in the estimated tree. This shows how small differences
in the sketch-based distance estimates can influence the final tree.


% ============================================================
\section{Conclusion}
% ============================================================

The program successfully estimates genomic distances using k-mer sketches
and uses these distances to construct binary evolutionary trees with UPGMA.

For Task 1, both sketching density and k-mer size affected the estimated
distance. The density experiment showed that increasing the density
increases the sketch size and changes the estimated distance. In the
validation experiment, estimates at higher densities tended to be closer
to the exact values, although the sketching error was not perfectly
monotonic.

The k-mer-size experiment showed a stronger effect. For G2--G3, the
exact distance increased from 0.4841 at $k=10$ to 0.8731 at $k=90$.

For Task 2, genomes with smaller estimated distances were generally
placed closer together in the UPGMA tree. The simulated lineage
comparison showed that Test 1 agreed with the simulated lineage, Test 2
showed partial agreement, and Test 3 grouped all three examined
same-lineage pairs directly. Thus, the estimated trees captured several
of the expected relationships, but the results were not perfect.
Sketching error can affect the estimated distance matrix and consequently
the final tree structure.

The runtime experiment also showed that the main computational cost came
from processing the k-mers and calculating pairwise distances. The
20-genome \textit{E. coli} dataset took approximately 7.5 minutes in
total, while the UPGMA calculation itself took only a small fraction of
a second.



\newpage
\section{References}

\begin{enumerate}

\item Biopython, \textit{Bio.Phylo.TreeConstruction module},
available at: \url{https://biopython.org/docs/dev/api/Bio.Phylo.TreeConstruction.html}

\item D. Cain, \textit{UPGMA: Efficient UPGMA implementation},
GitHub repository,
available at: \url{https://github.com/DavidCain/upgma}


\item
DA7067 Computational Mathematics,
\emph{MinHash--Jaccard lecture slides},
Stockholm University, 2026.

\end{enumerate}

\end{document}